% Configurazione
\documentclass{article}

\usepackage{titling} % Required for inserting the subtitle
\usepackage{graphicx} % Required for inserting images
\usepackage{tabularx} % Per l'ambiente tabularx (tabelle)
\usepackage{calc} % Sempre per le tabelle
\usepackage[hidelinks]{hyperref} % Per i collegamenti ipertestuali, ad esempio sulla table of contents
\usepackage[italian]{babel} % Per la lingua italiana nelle scritte automatiche
\usepackage{xcolor} % Per colorare il testo
\usepackage{lipsum} % Per generare lorem ipsum
\usepackage[normalem]{ulem} % Per sottolineare il testo
\usepackage{array} % Per la visualizzazione fluttuante di array di domande e risposte
\usepackage{ragged2e} % Pacchetto necessario per \justifying che giustifica il testo di tabelle

\newcommand{\ulhref}[2]{\href{#1}{\uline{#2}}} % Nuovo comando per sottolineare i link
\newcommand{\ulref}[1]{\uline{\ref{#1}}} % Nuovo comando per sottolineare i collegamenti a immagini

\graphicspath{ {immagini/} {../../../shared/immagini/} }


% Variabili
\newcommand{\data}{18 Ottobre 2024}


% Struttura
\begin{document}

% Parte costante
\begin{minipage}{0.4\textwidth}
    \includegraphics[width=0.6\textwidth]{Logo Universita di Padova.png}
\end{minipage}
\begin{minipage}{0.55\textwidth}
    \textcolor{red}{\textbf{Università degli Studi di Padova}} \\
    \textcolor{red}{Laurea: Informatica} \\
    \textcolor{red}{Corso: Ingegneria del Software} \\
    \textcolor{red}{Anno Accademico: 2024/2025}
\end{minipage}

\begin{minipage}{0.4\textwidth}
    \includegraphics[width=0.6\textwidth]{Logo SWEg.png}
\end{minipage}
\begin{minipage}{0.55\textwidth}
    \textbf{Gruppo: SWEg Labs} \\
    Email: \textsf{gruppo.sweg@gmail.com}
\end{minipage}

\vspace{2cm}
\centering\section*{\Huge \titolo}
\vspace{0.5cm}
\newpage
\tableofcontents
\newpage

% Parte variabile
\section{Informazioni generali}

\begin{itemize}
    \item \textbf{Tipo di riunione:} esterna con proponenti
    \item \textbf{Luogo:} meeting zoom
    \item \textbf{Data:} 22/10/2024
    \item \textbf{Ora inizio:} 17:00
    \item \textbf{Ora fine:} 17:40
    \item \textbf{Responsabile:} Stefani Riccardo
    \item \textbf{Notaio:} Verzotto Davide
    \item \textbf{Partecipanti:}
    \begin{itemize}
        \renewcommand{\labelitemii}{--}
        \item Bolognini Federica
        \item Fantinato Michael
        \item Loat Giacomo
        \item Righetto Filippo
        \item Stefani Riccardo
        \item Verzotto Davide
        \item Carlesso Gianluca (\emph{Ergon Informatica Srl})
        \item Tieppo Anna (\emph{Ergon Informatica Srl})
    \end{itemize}
\end{itemize}
\newpage
\section{Ordine del giorno}

\begin{itemize}
    \item Parlare con il proponente \emph{Ergon} e presentare alcune domande al fine di poter segnare le risposte rcevute, ed effettuare successivamente un'analisi più approfondita.
\end{itemize}

\newpage
\section{Diario della riunione}

\begin{itemize}
    \item Analizzata la presentazione del capitolato \emph{LLM: Assistente virtuale} di \emph{Ergon}, e raccolte domande e dubbi nel Foglio Google consultabile \ulhref{https://docs.google.com/spreadsheets/d/1W68ouX5T0ck2tDQm2uQCv7EYVG3xXCtpChtITe8GQg0/edit?gid=0}{qui}.
    
    \item Analizzata la documentazione del capitolato \emph{NearYou - Smart custom advertising platform} di \emph{SyncLab}, e abbiamo deciso di non procedere con l'interesse al capitolato, perchè abbiamo valutato l'impegno come aldilà delle nostre possibilità.
    
    \item Contattata via mail l'azienda \emph{Ergon} per organizzare un incontro online in cui rispondere e discutere delle domande presenti nei rispettivi Fogli Google.

    \item Accettato l'utilizzo di \emph{Overleaf} e \emph{GitHub Pages} per gestire la documentazione del gruppo

    \item Preparata una bozza di piano per l'approccio all'incontro di lunedì alle 17:00 con l'azienda \emph{Azzurro Digitale}. 
    \begin{itemize}
        \renewcommand{\labelitemii}{--}
        \item Inizialmente presentare qual è la nostra idea, cioè che cosa abbiamo capito del progetto, per dimostrare che lo abbiamo analizzato per bene.
        \item Presentare via Zoom il Foglio Google con le domande che abbiamo in relazione all'azienda (le quali saranno divise in "funzionali" e "non funzionali") e scrivere in live le risposte.
        \item Applicare il metodo MoSCoW come tecnica di prioritizzazione dei requisiti del progetto, cioè fare domanda all'azienda per capire cosa ne pensano di ciascun requisito.
    \end{itemize}
\end{itemize}
\newpage
\section{Todo}

Durante la riunione sono emersi i seguenti task da svolgere:

\vspace{0.5cm}

\begin{table}[htbp]
\begin{tabular}{|p{0.4\textwidth}|p{0.6\textwidth}|}
    \hline
    \multicolumn{1}{|c|}{\textbf{Assegnatario}} & \multicolumn{1}{|c|}{\textbf{Task Todo}} \\
    \hline
    Stefani Riccardo & Redigere verbale riunione \emph{17/10/24} \\
    \hline
\end{tabular}
\end{table}

\end{document}
